\chapter{Introdução}\label{chp:INTRODUCAO}

Na era da informação, a disseminação de novos meios de comunicação dentro do ambiente universitário, com o intuito de facilitar a comunicação entre discentes e docentes, vem se tornando cada vez mais relevante. Nesse cenário, os agentes conversacionais podem ser ferramentas importantes e ótimos intermediadores para a busca de informações de âmbito acadêmico, como respostas a FAQs relacionadas a matrículas, grade curricular e eventos, facilitando e otimizando o tempo das secretarias, que outrora possuíam o papel de entregar esse tipo de resposta.

Conforme já observado em estudos acadêmicos, tal como a monografia entitulada de \textit{``Solução Chatbot no Ambiente Acadêmico da UFRJ''} \cite{ufrj}, na qual é apresentada uma solução baseada em regras dentro do ambiente acadêmico da UFRJ, focada na entrega de informações sobre matrícula, calendário acadêmico, atestados e declarações, e também no trabalho \textit{``Um Chatbot Especializado para o Contexto da Universidade Federal de Ouro Preto''} \cite{ufop}, com aplicações similares e tecnologias mais recentes, observa-se que a utilização de chatbots pode ser aplicada com esses fins.

Atualmente, na Universidade Federal Rural do Rio de Janeiro, as atas de reuniões acadêmicas armazenam informações essenciais sobre o desenvolvimento do curso e decisões acadêmicas que afetam os alunos e os professores, representando um histórico valioso para preservação e consulta.

Nesse contexto, foi observada a necessidade de uma ferramenta que facilite a busca por atas que contenham informações pertinentes sobre assuntos específicos ou buscas mais genéricas. Com fins de pesquisa e levantamento do histórico de reuniões dentro do departamento e considerando que a quantidade de atas cresce continuamente, realizar a busca manual torna-se cada vez mais difícil. Além disso, a aplicação de um sistema de consulta de atas demanda manutenção, aprendizado por parte do operador e implantação adequada ao ambiente acadêmico.

\section{Objetivo}

O problema central deste trabalho está em levar discentes, docentes e o público acadêmico em geral até as atas que mais se adequam aos critérios de busca informados. Isso é feito por meio de um roteiro guiado de diálogo, com o intermédio de um chatbot baseado em regras para a comunicação e um sistema de busca dentro de uma base de dados composta por essas atas.

Este trabalho tem como objetivo disponibilizar uma ferramenta conversacional, na forma de um chatbot, que conduza discentes, docentes e o público acadêmico em geral por esse roteiro guiado de diálogo. O roteiro é estruturado em torno de um conjunto fechado de critérios de busca (como departamento, disciplina, docentes, discentes e período) e tem como fim encontrar as atas institucionais da UFRRJ mais pertinentes ao conteúdo da busca realizada.


\section{Organização do Trabalho}

O restante deste trabalho está organizado da seguinte forma: o Capítulo~\ref{chp:FUNDAMENTACAO} apresenta a fundamentação teórica necessária para a compreensão da proposta, abordando a origem histórica e a evolução dos chatbots, a definição formal do termo, a taxonomia dos tipos existentes e os fundamentos de Busca e Recuperação da Informação. O Capítulo~\ref{chp:PROPOSTA} descreve a proposta de solução, incluindo a arquitetura do sistema e a forma como seus componentes se relacionam para responder às consultas dos usuários. O Capítulo~\ref{chp:EXPERIMENTOS} apresenta a implementação da solução proposta, as ferramentas e tecnologias utilizadas em sua construção. Por fim, o Capítulo~\ref{chp:CONCLUSAO} traz as considerações finais do trabalho, suas limitações e sugestões para trabalhos futuros.
